\section{The original ordinary torus character and its full weight condition}
\label{sec:ordinary-theta-weight}

\status{OTW1--OTW12 reconstruct the complete original
\texttt{exists\_\allowbreak theta\_\allowbreak is\_\allowbreak weight\_\allowbreak of\_\allowbreak lam\_\allowbreak obj} endpoint.  The character
takes values in the original complete coefficient ring and is chosen
before every compact-open level.  The local reciprocity map is the
literal \texttt{artOf} certified by CIP4; the independent NTL17
construction supplies the original full \texttt{neatU}.
OTW13 gives the exact \texttt{UnitaryHeckeData} assembly once its
original integral \texttt{TypeForAlt} is supplied, using the already
completed CAP17 residual construction.  No type, ordinary module,
control theorem, arithmetic-fiber separation or common ideal is
produced by that conditional assembly.}

We use the pinned source throughout.  A complete local ring is
understood to be separated for its maximal-ideal topology.  Congruences
in a coefficient ring need not be congruences in a field: nilpotents
and nontrivial ramification are allowed.  The finite complex
characters used below detect finite abelian groups; they never change
the original coefficient ring.

\subsection{Complete-coefficient characters and the actual Artin map}

% BEGIN OTW1
\paragraph{OTW1. The original coefficients have finite adic quotients.}
\label{claim:OTW1}
Let $\ell$ be prime, let $O=c.O$ for the original
$c:\mathrm{Coef}(\ell)$, and let $R$ be an original
$\mathrm{IsCO}(O,R)$ algebra.  Then $\mathbb Z_\ell\to O\to R$
are local maps, $R$ is complete Noetherian local with finite residue
field, and every $R/\mathfrak m_R^N$ is finite.  In particular
these assertions hold for $R=\Lambda.\Lambda$ for every original
$\Lambda:\mathrm{LamObj}(r)$, and
\[
                 \bigcap_{N\geq0}\mathfrak m_R^N=0.
                 \tag{OTW.1}
\]
\begin{proof}
By definition $O$ is a DVR subalgebra of
$\overline{\mathbb Q}_\ell$, finite as a $\mathbb Z_\ell$-module.
The element $\ell$ is not a unit of $O$.  If it were, then
$O=\ell O$.  For finitely many $\mathbb Z_\ell$-generators
$a_1,\ldots,a_d$ we could write
$a_i=\ell\sum_j b_{ij}a_j$ with $b_{ij}\in\mathbb Z_\ell$.
Multiplication by the adjugate of $I-\ell B$ would give
$\det(I-\ell B)a_i=0$ for every $i$.  This determinant is congruent
to one modulo $\ell$, hence is a unit of $\mathbb Z_\ell$.
All generators would vanish, contradicting $1\ne0$ in $O$.
Thus $\ell\in\mathfrak m_O$.  The nonunits of
$\mathbb Z_\ell$ are exactly $\ell\mathbb Z_\ell$, and their
images lie in $\mathfrak m_O$; this proves locality.  Moreover
$O/\ell O$ is finite, since it is generated by the images of the
$a_i$ over $\mathbb F_\ell$.  Its quotient $O/\mathfrak m_O$
is therefore finite.

The original $\mathrm{IsCO}$ record already supplies that $R$ is
Noetherian local and complete, that $O\to R$ is local, and that
each $x\in R$ differs from an element of $O$ by an element of
$\mathfrak m_R$.  The induced map of residue fields is injective,
as a unital map of fields, and the last property makes it surjective.
Thus $R/\mathfrak m_R\simeq O/\mathfrak m_O$ is finite.
This verifies every field of the source's $\mathrm{IsCNL}(R)$.
Separatedness in the supplied adic completeness gives (OTW.1).

For the finite-quotient assertion, choose generators
$x_1,\ldots,x_e$ of $\mathfrak m_R$ by Noetherianity.
The finitely many monomials of degree $d$ in the $x_i$ span
$\mathfrak m_R^d/\mathfrak m_R^{d+1}$ over the finite residue
field.  This quotient is finite.  The exact sequences
\[
 0\longrightarrow\mathfrak m_R^d/\mathfrak m_R^{d+1}
 \longrightarrow R/\mathfrak m_R^{d+1}
 \longrightarrow R/\mathfrak m_R^d\longrightarrow0
 \tag{OTW.2}
\]
give finiteness by induction.  For $N=0$ the quotient is the
one-element ring, which is also finite.  The units of each quotient
form a finite abelian group, with no reducedness assumption.
Finally $\mathrm{LamObj}$ contains the original $\mathrm{IsCO}$
instance, so the stated specialization follows on that very ring.
\end{proof}
% END OTW1

% BEGIN OTW2
\paragraph{OTW2. The literal principal-unit subgroups form a neighborhood base.}
\label{claim:OTW2}
Let $k$ be a nonarchimedean local field, $\mathcal O_k$ its valuation
ring, and $\mathfrak p$ its maximal ideal.  The source's $U_o(k)$
is the image of $\mathcal O_k^\times$ in $k^\times$, and
\[
 U_c(k,d)=\{u\in\mathcal O_k^\times:u-1\in\mathfrak p^d\}
                  \subset k^\times.                 \tag{OTW.3}
\]
If $W\leq k^\times$ is open, there is one $c\geq0$ such that
\[
                d\geq c\quad\Longrightarrow\quad U_c(k,d)\leq W.
                \tag{OTW.4}
\]
\begin{proof}
A nonzero element has valuation one precisely when it and its inverse
are integral; this identifies $U_o(k)$ with the integral units.
The source defines $U_c(k,d)$ by mapping to $k^\times$ the kernel
of the unit reduction
$\mathcal O_k^\times\to(\mathcal O_k/\mathfrak p^d)^\times$.
The integral embedding is injective and reduction sends a unit to
one exactly when its difference from one belongs to $\mathfrak p^d$.
This proves (OTW.3), including $d=0$, when all integral units occur.

Choose a uniformizer $\pi$.  The sets $\pi^d\mathcal O_k$ form
a neighborhood base of zero in the local-field topology.  Inversion
is continuous at one, so $1+\pi^d\mathcal O_k$, for $d\geq1$,
form a neighborhood base of one in $k^\times$ and consist of
units.  Since $W$ is an open neighborhood of one, one such set is
contained in $W$.  If $d\geq c$, then
$\mathfrak p^d\subseteq\mathfrak p^c$, proving (OTW.4).
This argument concerns the actual local topology and congruence
subgroups; it does not require a separate general-linear bridge.
\end{proof}
% END OTW2

% BEGIN OTW3
\paragraph{OTW3. Finite abelian groups are separated by complex characters.}
\label{claim:OTW3}
For a finite abelian group $B$ and $b\ne1$ there is a homomorphism
$\phi_b:B\to\mathbb C^\times$ with $\phi_b(b)\ne1$.
\begin{proof}
If $b$ has order $r>1$, send it in its cyclic subgroup to a primitive
$r$th root of unity.  A character $\phi$ of any subgroup $H\leq B$
extends across one additional generator $x$.  Let $d>0$ be the least
integer with $x^d\in H$, and choose $z\in\mathbb C^\times$ with
$z^d=\phi(x^d)$.  Define
\[
                   \phi'(h x^j)=\phi(h)z^j.
                   \tag{OTW.5}
\]
If two expressions agree, then $x^{j-j'}\in H$, so
$j-j'=dq$ for an integer $q$.  Their two values agree because
$\phi(x^{dq})=z^{dq}$.  Thus (OTW.5) is well defined; commutativity
of $B$ gives multiplicativity, and it extends $\phi$.
Repeating this operation terminates because $B$ is finite.
The resulting character of $B$ retains the original nontrivial value
at $b$.  This proves the assertion without assuming a decomposition
of $B$ into cyclic factors.
\end{proof}
% END OTW3

% BEGIN OTW4
\paragraph{OTW4. Finite smooth characters are killed above one Artin depth.}
\label{claim:OTW4}
Fix the original number field $E$, algebraic closure $\Omega$,
valuation $A$ over a finite place $w$, and
$a:W_A\to k^\times$, $k=E_w$, satisfying the full original
$\mathrm{IsLocalArtin}$ predicate.  Let $f:W_A\to B$ be a
homomorphism to a finite abelian group which is smooth in the
source's sense: its kernel on the original inertia is open.
Then there is $c$ such that
\[
 d\geq c,\quad a(\sigma)\in U_c(k,d)
                 \quad\Longrightarrow\quad f(\sigma)=1.
                 \tag{OTW.6}
\]
\begin{proof}
For every $b\in B\setminus\{1\}$ choose the character of OTW3.
The inertia kernel of $\phi_b f$ contains that of $f$, hence is
open: a subgroup containing an open subgroup is a union of its open
cosets.  The exact \texttt{factor} field of $\mathrm{IsLocalArtin}$
therefore gives a character $\xi_b:k^\times\to\mathbb C^\times$
with open kernel and
\[
                        \phi_b f=\xi_b a.
                        \tag{OTW.7}
\]
The finite intersection $W=\bigcap_{b\ne1}\ker\xi_b$ is open.
Choose $c$ by OTW2.  Under the hypotheses of (OTW.6), every
$\phi_b(f(\sigma))$ equals one.  If $f(\sigma)\ne1$, choosing
$b=f(\sigma)$ contradicts OTW3.  The case $B=\{1\}$ is
immediate and also follows using the empty intersection.

Only the literal smooth-character factor property and the local
congruence topology enter this argument.  In particular it does not
assume that the source Weil subgroup has a different, unstated
topology or that $f$ has open kernel in the entire decomposition group.
\end{proof}
% END OTW4

% BEGIN OTW5
\paragraph{OTW5. Descent of a character valued in the whole complete ring.}
\label{claim:OTW5}
Retain $a$ of OTW4, and let $R$ be any $\mathrm{IsCNL}$ ring with
maximal ideal $\mathfrak m$.  Suppose
$\psi:D_A\to R^\times$ satisfies the original
$\mathrm{AdicCont}(\mathfrak m,\psi)$, meaning that its underlying
ring-valued function is locally constant modulo every
$\mathfrak m^N$.  There is a unique homomorphism
$\chi:k^\times\to R^\times$ such that
\[
 \chi(a(x))=\psi(g)
 \quad\hbox{if }x\in W_A,\ g\in D_A
       \hbox{ have the same underlying Galois automorphism}.
 \tag{OTW.8}
\]
For this same $\chi$,
\[
 \forall N\ \exists c\ \forall d\geq c\ \forall y\in U_c(k,d),
                       \chi(y)-1\in\mathfrak m^N.
                       \tag{OTW.9}
\]
\begin{proof}
Restrict $\psi$ to the original Weil subgroup and call the restriction
$f$.  OTW1's finite-quotient argument applies to any $\mathrm{IsCNL}$
ring.  Thus
\[
 f_N:W_A\longrightarrow(R/\mathfrak m^N)^\times
 \tag{OTW.10}
\]
has finite abelian target.  It is smooth: the inclusion of the
original inertia into $D_A$ is continuous, and local constancy of
$\psi$ modulo $\mathfrak m^N$ makes the inverse image of one on
inertia open.  Equality of units in a quotient is equality of their
underlying ring elements.  OTW4 gives $c_N$ with
\[
 d\geq c_N,\quad a(x)\in U_c(k,d)
               \quad\Longrightarrow\quad f(x)-1\in\mathfrak m^N.
               \tag{OTW.11}
\]
If $a(x)=1$, use $d=c_N$ in (OTW.11) for every $N$.
Separatedness gives $f(x)=1$.  Consequently
$\ker a\subseteq\ker f$ in the actual groups, before taking
any coefficient quotient.

The \texttt{surj} field of the same $\mathrm{IsLocalArtin}$ says
that $a$ is onto.  Set $\chi(y)=f(x)$ for any $x$ with $a(x)=y$.
Two choices differ by an element of $\ker a$, which $f$ kills;
the value is independent of the choice.  Lifting two elements and
their product proves multiplicativity.  Surjectivity also proves
uniqueness.  If $g$ and $x$ have the same underlying automorphism,
the decomposition-group element $g$ is precisely the original
inclusion of $x$, which proves (OTW.8).

For (OTW.9), keep the previously chosen $c_N$.  For
$y\in U_c(k,d)$ choose $x$ with $a(x)=y$ and apply (OTW.11).
This proves the asserted congruence for the single full-ring
character just constructed.  We have not chosen incompatible
characters separately at different precisions.  Completeness is
part of the source coefficient category; separatedness is the part
of it used in the kernel argument.  The proof works for nonreduced
coefficient rings as well.
\end{proof}
% END OTW5

% BEGIN OTW6
\paragraph{OTW6. Cyclotomic and multiplicative adic continuity.}
\label{claim:OTW6}
Let $j:\mathbb Z_\ell\to R$ be a local ring homomorphism and
$e:\Gamma\to\mathbb Z_\ell$ a continuous function.  Then
$j e$ is adically continuous for $\mathfrak m_R$.  Adically
continuous unit-valued characters are closed under multiplication,
nonnegative integral powers and inversion.  For the original
$H=\mathrm{hset}(Q,\tau_0)$ and $R=\Lambda.\Lambda$, this applies
to the literal $H.\mathrm{epsV}(R,v)$ at every $v$.
\begin{proof}
For every $N$ and $x\in\Gamma$, continuity gives a neighborhood
on which $e(y)-e(x)\in\ell^N\mathbb Z_\ell$.
Since $j$ is local, $j(\ell)\in\mathfrak m_R$, so
\[
              j(e(y))-j(e(x))\in\mathfrak m_R^N.
              \tag{OTW.12}
\]
This is exactly local constancy after reduction modulo
$\mathfrak m_R^N$, including the trivial $N=0$ quotient.
Products and powers descend to the corresponding operations in
each quotient.  For units, if $u-u'\in\mathfrak m_R^N$ then
\[
              u^{-1}-(u')^{-1}=-u^{-1}(u')^{-1}(u-u')
                         \in\mathfrak m_R^N.          \tag{OTW.13}
\]
Intersecting the finitely many neighborhoods proves all the
closure assertions.  No topology on the entire ring of abstract
functions, or continuity of inversion on such a ring, is needed.

For the actual cyclotomic character, its reduction modulo
$\ell^N$ is the action on the finitely many $\ell^N$th roots of
unity.  A finite extension contains these roots, and its pointwise
stabilizer is open in the original Krull topology.  Hence each
reduction is locally constant and the resulting
$\mathbb Z_\ell^\times$-valued character is continuous.  Restricting
to the original $D_{s.A(v)}$ preserves continuity.
The definition of $\mathrm{epsV}$ is exactly this restriction
followed by
\[
             j=(O\longrightarrow R)\circ(\mathbb Z_\ell\longrightarrow O).
             \tag{OTW.14}
\]
OTW1 makes $j$ local.  Formula (OTW.12) therefore proves the
required adic continuity for the actual $\mathrm{epsV}$, with no
renormalization of the cyclotomic character.
\end{proof}
% END OTW6

% BEGIN OTW7
\paragraph{OTW7. The exact inverse cyclotomic adjustment.}
\label{claim:OTW7}
Let $\psi,\varepsilon:D_A\to R^\times$ be adically continuous,
with $R$ as in OTW5, and let $i\geq0$.  There exists one
$\chi_i:k^\times\to R^\times$ satisfying (OTW.9) and
\[
        \psi(g)\,\varepsilon(g)^i\,\chi_i(a(x))=1
        \quad\hbox{whenever }g\hbox{ and }x
                 \hbox{ have the same underlying automorphism}.
        \tag{OTW.15}
\]
\begin{proof}
Apply OTW5 to the character
\[
                     \eta=(\psi\,\varepsilon^i)^{-1}.
                     \tag{OTW.16}
\]
It is a homomorphism because $R^\times$ is abelian, and it is
adically continuous by OTW6.  The resulting $\chi_i$ obeys
$\chi_i(a(x))=(\psi(g)\varepsilon(g)^i)^{-1}$, which gives
(OTW.15) in the displayed order.  Its congruence continuity is
(OTW.9).

This inverse is essential.  In the later original split CM
application, CIP4 fixes the actual Artin convention
$\operatorname{ord}(a(x))=-\deg(x)$, and CIP5 identifies the
cyclotomic character on inertia with the positive local norm of
$a(x)$.  We do not replace $a$ by an arithmetic-normalized map,
nor replace (OTW.16) by $\psi^{-1}\varepsilon^i$.
No division by $\ell$, by $n!$, or by a difference of ordered
characters occurs.
\end{proof}
% END OTW7

\subsection{The actual Iwahori torus and its uniform congruence depth}

% BEGIN OTW8
\paragraph{OTW8. Diagonal coordinates on the original local torus.}
\label{claim:OTW8}
For the source's $\mathrm{glIwa}(k,n)$ there is a homomorphism
\[
       d:T\longrightarrow(k^\times)^n,
       \qquad \operatorname{diag}(d(t))=t,             \tag{OTW.17}
\]
whose image is $(\mathcal O_k^\times)^n$.  Its coordinates are
unique, it sends $e_T(i,y)$ to the tuple with $y$ in position $i$
and one elsewhere, and
\[
 t\in T_c(c)\quad\Longleftrightarrow\quad
                d(t)_i\in U_c(k,c+1)\quad(0\leq i<n).
                \tag{OTW.18}
\]
\begin{proof}
The original \texttt{glData.T} is the image of the integral unit
torus under the diagonal map, and \texttt{glIwa} retains that
underlying subgroup.  A diagonal matrix uniquely determines all
of its diagonal entries.  Taking those entries defines $d$;
diagonal multiplication is coordinatewise, so $d$ is a homomorphism.
OTW2 identifies its coordinates with integral units.  Conversely,
every tuple of integral units gives a diagonal element of $T$.
The definition of $e_T$ is exactly the one-coordinate diagonal,
which proves its asserted image.

The original $T_c(c)$ is $T\cap K_c(k,n,c+1)$.  An element of
$K_c(k,n,c+1)$ is the image of an invertible integral matrix
congruent to the identity modulo $\mathfrak p^{c+1}$.
If that image is the diagonal matrix $t$, injectivity of
$\mathcal O_k\to k$ forces the integral matrix to be the diagonal
matrix with entries $d(t)_i$.  Its reduction is the identity
exactly when $d(t)_i-1\in\mathfrak p^{c+1}$ for every $i$.
Conversely the diagonal integral lift of such a tuple is invertible
over $\mathcal O_k$ and reduces to the identity.  OTW2 gives
(OTW.18).  The index is precisely $c+1$, as in the source, not $c$.
The assertions also hold for $n=0$, with the empty tuple and the
one-element torus.
\end{proof}
% END OTW8

% BEGIN OTW9
\paragraph{OTW9. One character on the product torus, with one depth per precision.}
\label{claim:OTW9}
Let $\mathcal V$ be finite and $k_v$ the original local fields.
Suppose characters $\chi_{v,i}:k_v^\times\to R^\times$ are given
for $v\in\mathcal V$ and $i\in\operatorname{Fin}(n)$, and
\[
 \forall v,i,N\ \exists c_{v,i,N}\ \forall d\geq c_{v,i,N}
 \ \forall y\in U_c(k_v,d),\quad
                       \chi_{v,i}(y)-1\in\mathfrak m_R^N.
                       \tag{OTW.19}
\]
On the original $I=\mathrm{iwS}(n,S)$ define
\[
        \theta(t)=\prod_{v\in\mathcal V}\prod_{i=0}^{n-1}
                         \chi_{v,i}(d_v(t_v)_i).
                         \tag{OTW.20}
\]
Then $\theta:I.T\to R^\times$ is a homomorphism,
\[
                  \theta(e_V(v,i,y))=\chi_{v,i}(y),
                  \tag{OTW.21}
\]
and it satisfies the entire original $I.\mathrm{ContChar}$:
\[
 \forall N\ \exists c\ \forall t\in I.T,
             t\in I.T_c(c)\ \Longrightarrow\ \theta(t)-1\in\mathfrak m_R^N.
             \tag{OTW.22}
\]
\begin{proof}
The definitions of \texttt{piData}, \texttt{piIwa} and
\texttt{iwS} give $I.T=\prod_vT_v$ and
$I.T_c(c)=\prod_v(T_v)_c(c)$.  The source projections are
the restrictions of the original coordinate evaluations.
All the factors in (OTW.20) are homomorphisms to the abelian
group $R^\times$, so their finite product is a homomorphism.
The source $e_V$ is the one-place, one-coordinate insertion of
$e_T$.  OTW8 makes every factor except $(v,i)$ equal to one;
the remaining factor is $\chi_{v,i}(y)$.  This proves (OTW.21).

For fixed $N$, choose the finitely many depths in (OTW.19) and put
\[
                   c=\max_{(v,i)}c_{v,i,N},            \tag{OTW.23}
\]
with value zero for an empty index set.  If $t\in I.T_c(c)$,
OTW8 puts $d_v(t_v)_i$ in $U_c(k_v,c+1)$, and
$c+1\geq c_{v,i,N}$.  Each factor of (OTW.20) is therefore
one modulo $\mathfrak m_R^N$.  Reducing the finite product in
the quotient ring proves (OTW.22).
The character was defined before $N,c,t$; only its congruence
depth depends on $N$.  No choice of a compact-open level appears
in (OTW.20) or (OTW.23).
\end{proof}
% END OTW9

% BEGIN OTW10
\paragraph{OTW10. The literal \texttt{IsTheta} character on the original \texttt{hset}.}
\label{claim:OTW10}
Retain every original parameter of the source endpoint:
$\ell,L,\Omega$, a $\mathrm{HidaSetting}$ $s$, $n\geq1$,
$V:s.\mathrm{Lev}$, $Q:\mathrm{QData}(V)$, $c:\mathrm{Coef}(\ell)$,
any integral representation
\[
 \tau_0:\mathrm{KT}(n,\mathrm{wt}(s))
                      \longrightarrow\mathrm{GL}_m(c.O),
 \tag{OTW.24}
\]
an original $r:\mathrm{Resid}(V.\mathrm{frame},c.O,n)$ and
$\Lambda:\mathrm{LamObj}(r)$.  Put $H=\mathrm{hset}(Q,\tau_0)$
and $R=\Lambda.\Lambda$.  There is one homomorphism
$\theta:H.X.I.T\to R^\times$ with
$\mathrm{IsTheta}(Q,\tau_0,\Lambda,\theta)$ and
$H.X.I.\mathrm{ContChar}(\theta)$.  Moreover $\mathrm{IsTheta}$
alone uniquely determines this character on the original torus.
\begin{proof}
The finite index set $\mathcal V=V.\mathrm{frame}.V_l$ is exactly
the set of places of $L^+$ above $\ell$.  The original frame sets
\[
 D_v=\operatorname{Dec}_{L}(s.A(v)),\qquad
 I_v=\operatorname{Ine}_{L}(s.A(v))\ \text{as a subgroup of }D_v,
 \qquad k_v=L_{s.w(v)}.                              \tag{OTW.25}
\]
The $\mathrm{hset}$ uses these same places, valuations and
$a_v=\mathrm{artOf}(s.hw(v))$.  CIP4 proves that each actual
$a_v$ satisfies the full $\mathrm{IsLocalArtin}$ predicate.
Its hypotheses are met: $L$ is the given CM field, $s.A(v)$ is
the given place over $s.w(v)$, and $s.\mathrm{split}$ applies
at every $v\mid\ell$, so $s.w(v)$ is distinct from its conjugate.
No new local or global reciprocity map is selected here.

For each ordered pair $(v,i)$ the \texttt{inert} field of
$\Lambda$ gives a character
$\psi_{v,i}:D_v\to R^\times$ satisfying
$r.\mathrm{IsPsi}(v,i,\psi_{v,i})$ and
\[
                         \kappa_{v,i}=\psi_{v,i}|_{I_v}.
                         \tag{OTW.26}
\]
Its \texttt{cont} field is exactly the adic continuity used in OTW5.
OTW1 supplies $\mathrm{IsCNL}(R)$, and OTW6 supplies adic
continuity of the actual $\varepsilon_v=H.\mathrm{epsV}(R,v)$.
Use OTW7 for this $\psi_{v,i}$, this $\varepsilon_v$ and the
original nonnegative integer $i$, and then OTW9.  All choices
have been made on the fixed original data before any later level.

Let $x$ be an original Weil element, let $g\in D_v$ have the
same underlying Galois automorphism, suppose $g\in I_v$, and let
$y\in U_o(k_v)$ satisfy $a_v(x)=y$.  Equations
(OTW.15), (OTW.21) and (OTW.26) give, in precisely the source order,
\[
       \kappa_{v,i}(g)\,\varepsilon_v(g)^i\,
                       \theta(e_V(v,i,y))=1.          \tag{OTW.27}
\]
This verifies every quantified variable and guard in
\texttt{IsTheta}; the finite-product continuity is (OTW.22).
In particular $\tau_0$ need not satisfy any type condition for
this construction.

For uniqueness, CIP4 also proves that $a_v$ maps original inertia
onto $\mathcal O_{k_v}^\times$: a Weil preimage of a unit has
degree zero by the same \texttt{ord\_eq}, hence lies in inertia.
Thus for each $v,i,y$, (OTW.27) forces the value of any proposed
$\theta$ on $e_V(v,i,y)$.  These elements generate the actual
product torus: a diagonal tuple is the finite product of its
one-place, one-coordinate diagonal factors.  Hence two characters
satisfying $\mathrm{IsTheta}$ agree everywhere.  The torus
character is therefore independent of the auxiliary extensions
$\psi_{v,i}$ chosen from \texttt{inert}; no compatibility of
independent finite-precision choices is being assumed.
\end{proof}
% END OTW10

\subsection{Every compact-open level and the exact data assembly}

% BEGIN OTW11
\paragraph{OTW11. The original neat stabilizer condition implies full freeness.}
\label{claim:OTW11}
For any original frame $X$, subgroup $U\leq X.U_b$, commutative
coefficient ring $A$, representation
$\tau:X.U_b\to\mathrm{GL}_m(A)$ and character
$\theta:X.I.T\to A^\times$, one has
\[
                   X.\mathrm{Neat}(U) \Longrightarrow\
                          X.\mathrm{Free}(\tau,\theta,U).
                   \tag{OTW.28}
\]
\begin{proof}
The source \texttt{Neat} says that for every $g,\gamma,u,b$
with $u\in U$, $b\in X.I.\mathrm{Iw}(0)$ and
\[
                     X.\mathrm{glob}(\gamma)g=g(u,b),
                     \tag{OTW.29}
\]
one has $u=1$ and $b=1$.  This makes $b$ central and puts $b$
in $X.I.T$, as required by the first two parts of \texttt{Free}.
Using that exact membership witness, its remaining matrix equality is
\[
                       \theta(1)\,\tau(1)=1\cdot I_m=I_m.
                       \tag{OTW.30}
\]
Both identities follow from the homomorphism laws.  This proves
the entire original \texttt{Free}, with no smoothness or type
condition on $\tau$, no continuity requirement on $\theta$ and
no condition on $m$.  It also applies after the original scalar
extension of $\tau$ to a complete coefficient ring.
\end{proof}
% END OTW11

% BEGIN OTW12
\paragraph{OTW12. The whole original theta-and-weight producer.}
\label{claim:OTW12}
For all the original data in OTW10, and in particular under the
source's stated \texttt{LiteratureRankNHypothesis}, there is one
\[
        \theta:(\mathrm{hset}(Q,\tau_0)).X.I.T\to(\Lambda.\Lambda)^\times
        \tag{OTW.31}
\]
which satisfies both original predicates
\[
        \mathrm{IsTheta}(Q,\tau_0,\Lambda,\theta),
        \qquad (\mathrm{hset}(Q,\tau_0)).\mathrm{IsWt}(\theta).
        \tag{OTW.32}
\]
The quantifier order is one $\theta$ before \emph{every} original
compact-open level $K$, with $Q$, $\tau_0$, $r$ and $\Lambda$ fixed.
\begin{proof}
Choose $\theta$ once by OTW10.  It already satisfies
$\mathrm{IsTheta}$ and the \texttt{cont} field of \texttt{IsWt}.
The remaining field is literally
\[
 \forall K:H.\mathrm{Lev},\quad
           H.X.\mathrm{Free}(H.\tau_A(R),\theta,H.\mathrm{lev}(K)),
 \quad H=\mathrm{hset}(Q,\tau_0),\quad R=\Lambda.\Lambda.
 \tag{OTW.33}
\]
Here $H.\mathrm{Lev}$ consists of the choices of a compact open
subgroup at each original added $Q$-place, and
\[
             H.\mathrm{lev}(K)=H.U_Q\vee
                       \bigvee_{q\in Q.Q}H.\mathrm{sec}(w_q)(K_q).
             \tag{OTW.34}
\]
The independent full neat-level theorem NTL17 gives
$H.X.\mathrm{Neat}(H.\mathrm{lev}(K))$ for each such $K$ on
these exact $\ell,L,\Omega,s,n,V,Q,c,m,\tau_0$.
It uses $n\geq1$ and the original Hida setting, including its
fixed auxiliary place, and imposes no \texttt{TypeForAlt} or
$H.\mathrm{Ax}$ premise.  Apply OTW11 with the actual
$H.\tau_A(R)$, the structural scalar extension of
$\tau_0\circ\mathrm{toKT}$, and with the already fixed $\theta$.
This proves (OTW.33) for every $K$ and completes \texttt{IsWt}.

The source uses its literature predicate to produce a local Artin
map.  Our inherited CIP4 supplies precisely its chosen \texttt{artOf}
on these original split CM places, with the checked class field
theory application already recorded there.  Thus no broader local
Langlands, integral-type or ordinary-control input has entered this
character construction.  All original coefficient primes allowed
by $s$ are included; in particular no additional exclusion of
$\ell=3$ or rank-four restriction is imposed.
\end{proof}
% END OTW12

% BEGIN OTW13
\paragraph{OTW13. Exact \texttt{HdEx} assembly after the original type is supplied.}
\label{claim:OTW13}
Fix the same original parameters as above and the original
$\iota:\overline{\mathbb Q}_\ell\simeq\mathbb C$.
Suppose an integer $m$ and integral $\tau_0$ as in (OTW.24)
are supplied with the complete original
$\mathrm{TypeForAlt}(r,\iota,\tau_0)$.  Then the original
$\mathrm{UnitaryHeckeData}(Q,\iota,\Lambda)$ is nonempty.
\begin{proof}
Apply OTW12 to this same $m,\tau_0,r,\Lambda$ to obtain
$\theta$ with its \texttt{theta} and \texttt{isWt} fields.
CAP17 applies to every original $V,Q,c,m,\tau_0,r$ and constructs
one function
\[
                      a:H.\mathrm{Op}\longrightarrow c.O
                      \tag{OTW.35}
\]
with the full $\mathrm{IsRes}(Q,r,\tau_0,a)$.
Its Frobenius field uses $V.\mathrm{pr}(\sigma^{-1})$ and the
original normalized Hecke polynomial; its uniformizer field uses
the ordered product $\prod_{i\leq j}\bar\psi_{v,i}\varepsilon_v^i$.
The validity and exact convention of \texttt{artOf} are the same
CIP4 input used here, so the two constructions do not choose
different reciprocity maps.

Take the record with the eight entries
\[
        (m,\tau_0,\mathrm{TypeForAlt},\theta,
                       \mathrm{IsTheta},\mathrm{IsWt},a,\mathrm{IsRes}).
        \tag{OTW.36}
\]
These are exactly the fields of the original data structure.
No transport to another coefficient field, level datum, type,
residual representation or weight algebra is needed.
The chosen $\theta$ and $a$ are held fixed when forming
$H.M_K(\theta,a,K)$ for every later $K$.

This proves the assembly performed by \texttt{HdEx} on the
stated type antecedent.  The original unconditional supplier of
that antecedent remains separate: \texttt{TypeForAlt} quantifies
over every original Teichm\"uller lift and every allowed
\texttt{IsRecT}, and tests the prescribed smooth irreducible
local types.  It has not been replaced by bare existence of an
integral representation.  Consequently OTW13 alone does not claim
the whole original \texttt{nonempty\_unitaryHeckeData} theorem
from \texttt{HypothesesH'} without that type producer.
\end{proof}
% END OTW13

\paragraph{Application boundary.}
The new completed producer is the entire original \texttt{TThetaEx},
using the independently reconstructed full \texttt{neatU}; its output
is not conditional on a newly assumed character or freeness theorem.
The exact residual \texttt{IsRes} producer is inherited from CAP17.
The original integral type, the remaining $H.\mathrm{Ax}$ fields,
ordinary module and control, Hecke determinant and separating
arithmetic fibers, the four Hida endpoints and all-level
\texttt{HasAut} on the same ideal remain outside this packet.
The selected Count-II whole leaf is therefore not declared closed.
All source/proof bindings concern the human reconstruction; no new
Lean replay is claimed.
